% Curva de luz de un tránsito planetario — el patrón que hay que aprender a leer.
%
% Diagrama propio en TikZ (sin licencias de terceros, vectorial y editable). Se
% tiñe solo con el color de la línea (milcGradeDark). Muestra el brillo de una
% estrella cayendo un poquito cuando un planeta pasa por delante, y volviendo:
% ese bajón simétrico es la firma de un exoplaneta.
%
% Se incrusta vía \guiaDiagrama desde el bloque `recursos:` del YAML.
\begin{tikzpicture}[scale=1,font=\sffamily]

  % ── Ejes ─────────────────────────────────────────────────────────────
  \draw[->,milcNegro!70] (0,0) -- (11,0) node[right,font=\scriptsize\bfseries] {tiempo (horas)};
  \draw[->,milcNegro!70] (0,0) -- (0,3.6) node[above,font=\scriptsize\bfseries,align=center] {brillo};

  % Referencia de brillo normal (100 %)
  \draw[dashed,milcNegro!30] (0,3) -- (11,3);
  \node[font=\scriptsize,milcNegro!55,left] at (0,3) {100\%};
  \node[font=\scriptsize,milcNegro!55,left] at (0,1.6) {99\%};

  % ── Curva: brillo estable → bajón simétrico → brillo estable ─────────
  \draw[line width=1.3pt,milcGradeDark]
    (0.3,3) -- (3.2,3)
    .. controls (3.7,3) and (3.8,1.6) .. (4.4,1.6)   % entrada del planeta
    -- (6.6,1.6)                                       % planeta cruzando (fondo plano)
    .. controls (7.2,1.6) and (7.3,3) .. (7.8,3)       % salida del planeta
    -- (10.6,3);

  % ── Puntos de datos (lo que el estudiante realmente mide) ────────────
  \foreach \x/\y in {0.6/3,1.4/3,2.2/3,3.0/3,4.6/1.6,5.5/1.6,6.4/1.6,8.2/3,9.2/3,10.1/3} {
    \fill[milcMostaza,draw=milcNegro!50] (\x,\y) circle (2pt);
  }

  % ── Anotaciones ──────────────────────────────────────────────────────
  \draw[<->,milcTurquesa,line width=.8pt] (4.4,1.15) -- (6.6,1.15);
  \node[font=\scriptsize\bfseries,milcTurquesa,below] at (5.5,1.15) {el planeta cruza};

  \draw[->,milcGradeDark,line width=.7pt] (2.0,2.35) -- (2.0,2.9);
  \node[font=\scriptsize\bfseries,milcGradeDark,align=center] at (2.0,2.1) {estrella\\sola};

  \node[font=\scriptsize\bfseries,milcNegro!70,align=center] at (5.5,3.35) {el bajón};

  % (La síntesis va en el pie de figura vía \guiaDiagrama, para no repetirla
  %  dentro del propio diagrama.)

\end{tikzpicture}
